LLM juries, ensemble labeling, and synthetic-data filters all rely on the same assumption: when several independent language models agree on an answer, the answer is probably correct.
We test how far this assumption holds and whether it is weaker for ethical judgement, where models may share biases.
We run 8 open-weight instruction models from 7 developer families on 3,500 items from 7 tasks: factual QA, math, two \ethics subsets, and a binary factual control matched to the ethics answer format.
The agreement rate (the fraction of models giving the plurality answer) is a strong label-free correctness signal on factual and math tasks (AUROC 0.78--0.91, point-biserial $r$ 0.37--0.61).
It is clearly weaker on ethics (AUROC 0.63--0.71, $r$ 0.20--0.33; pooled factual minus ethics $\Delta$AUROC $=+0.12$, 95\% CI $[0.08, 0.17]$).
However, the binary factual control reaches only AUROC 0.735, which is indistinguishable from \ethics commonsense, so answer format explains about half of the gap.
The plurality vote beats the best single model only on math ($+2.0$ to $+2.4$\pp). It is 1.8--2.4\pp below the best model on factual QA, a difference that is not significant, and significantly below it on ethics ($-4.8$\pp, Holm $p=0.008$).
On ethics, model errors are more correlated than on the control (error $\phi$ 0.45--0.49 vs.\ 0.30), and the vote's deficit grows as models are added.
Cross-model agreement is therefore a useful triage signal for factual and math answers, but it should not replace human labels for moral judgements: 10--17\% of unanimous \ethics verdicts contradict the human label.
